% =============================================================
% lx.tex  —  Plain TeX macro package
%
% Plain TeX parallel to lx.sty. Public API is identical where
% possible; divergences from LaTeX are noted explicitly.
%
% USAGE:
%   \input lx
%   \lxSetupPage             % call explicitly after setting flags/dimensions
%                            % (same in both Plain TeX and LaTeX)
%
% FLAGS (call these anywhere before \lxSetupPage):
%   \lxBindingLeft      \lxBindingRight
%   \lxDoubleSided      \lxSingleSided
%   \lxNotesOuter       \lxNotesInner
%
%   NOTE: in Plain TeX there is no documentclass option to agree
%   with; \lxDoubleSided / \lxSingleSided take full effect freely.
%   (In LaTeX, \lxDoubleSided must agree with twoside/oneside option,
%   which lx.sty cannot change after the fact.)
%
% MARGIN WIDTH — two mutually exclusive commands:
%   \lxMarginWidth=3truecm            % absolute (default: 3truecm)
%   \def\lxMarginWidthFraction{0.2}   % relative: fraction of page width
%
% MARGIN SEP:
%   \lxMarginSep=0.5em  (default)
%
% NOTE FONT:
%   \def\lxNoteFont{\the\scriptfont0 \spacing{7}{10}}  (default: complete 7pt environment)
%   Override with your own size-switching macro, e.g.: \def\lxNoteFont{\eightpoint}
%   (In LaTeX: \renewcommand{\lxNoteFont}{\footnotesize} — approximate)
%
% SIDE / TOP / BOTTOM MARGINS:
%   \lxInnerMargin, \lxOuterMargin, \lxTopMargin, \lxBottomMargin
%   (defaults: 1in each, matching Plain TeX's own defaults)
%
% HEAD / FOOT HOOKS:
%   \lxHeadInner  \lxHeadMiddle  \lxHeadOuter
%   \lxFootInner  \lxFootMiddle  \lxFootOuter
%
%   Defined as \newtoks registers (Plain TeX idiom) rather than
%   \newcommand (LaTeX idiom) — the only unavoidable API divergence.
%   Assign with: \lxHeadInner={...}
%
%   Defaults reproduce Plain TeX's own \headline and \footline:
%     head: all empty (LaTeX's default also has no running head)
%     foot: \lxFootMiddle={\the\pageno} (centred page number)
%
%   In double-sided mode inner/outer alternate per page via \ifodd\pageno.
%   In single-sided mode inner=left, outer=right throughout.
%
% MARGIN NOTES:
%   \lxMarginNote{NOTE TEXT}
%     Auto-incrementing form (default). Increments \lxNoteNo and
%     uses \lxNoteRef as the reference symbol (default: arabic integer).
%
%   \lxMarginNoteRef{REF}{NOTE TEXT}
%     Explicit-reference form. REF is any user-chosen symbol or string
%     (★, †, $\star$, 42, …); does not increment \lxNoteNo.
%     Allowable REF characters: any token(s) valid as a math superscript,
%     except unescaped % (comment), # (parameter), or
%     unbalanced braces.
%
%   \lxNoteNo  — counter (analogous to Plain's \pageno)
%   \lxNoteRef — formatter macro; override to change symbol style
%
% PAGE LAYOUT:
%   \lxSetupPage — call once after setting all flags and dimensions.
%   Must be called explicitly in both Plain TeX and LaTeX.
%   (lx.sty no longer uses \AtBeginDocument — see lx.sty PAGE SETUP.)
% =============================================================

% Guard against double-loading:
\expandafter\ifx\csname lxLoaded\endcsname\relax \else \endinput \fi
\expandafter\let\csname lxLoaded\endcsname=\relax
%                                         (no \NeedsTeXFormat in Plain TeX)
%                                         (no \ProvidesPackage in Plain TeX)
%
%                                         (no \RequirePackage in Plain TeX)
%                                         (packages are implicit in plain.tex)
%                                         (geometry, fancyhdr etc. not used)
%
% =============================================================
% SPACING UTILITY
%
% \spacing{NUM}{DEN} scales all skip parameters by the fraction
% NUM/DEN using integer arithmetic, changing the complete
% typographic environment (baselineskip, lineskip, parskip,
% display skips, jot etc.) consistently in one call.
% Usage: \spacing{7}{10}  — scale from 10pt environment to 7pt
%        \spacing{8}{10}  — scale from 10pt environment to 8pt
% (No LaTeX equivalent needed — \footnotesize etc. handle this.)
% =============================================================

\def\spacing#1#2{%
  \multiply\smallskipamount          by#1 \divide\smallskipamount          by#2
  \multiply\medskipamount            by#1 \divide\medskipamount            by#2
  \multiply\bigskipamount            by#1 \divide\bigskipamount            by#2
  \multiply\topskip                  by#1 \divide\topskip                  by#2
  \multiply\splittopskip             by#1 \divide\splittopskip             by#2
  \multiply\lineskip                 by#1 \divide\lineskip                 by#2
  \multiply\baselineskip             by#1 \divide\baselineskip             by#2
  \multiply\lineskiplimit            by#1 \divide\lineskiplimit            by#2
  \multiply\parskip                  by#1 \divide\parskip                  by#2
  \multiply\abovedisplayskip         by#1 \divide\abovedisplayskip         by#2
  \multiply\abovedisplayshortskip    by#1 \divide\abovedisplayshortskip    by#2
  \multiply\belowdisplayskip         by#1 \divide\belowdisplayskip         by#2
  \multiply\belowdisplayshortskip    by#1 \divide\belowdisplayshortskip    by#2
  \multiply\jot                      by#1 \divide\jot                      by#2}

% =============================================================
% FLAGS
% =============================================================

\newif\iflxBindingLeft  \lxBindingLefttrue   % default: binding=left
\newif\iflxDoubleSided  \lxDoubleSidedtrue   % default: sides=double
\newif\iflxNotesOuter   \lxNotesOutertrue    % default: notes=outer

\def\lxBindingLeft{\lxBindingLefttrue}
\def\lxBindingRight{\lxBindingLeftfalse}
\def\lxDoubleSided{\lxDoubleSidedtrue}
\def\lxSingleSided{\lxDoubleSidedfalse}
\def\lxNotesOuter{\lxNotesOutertrue}
\def\lxNotesInner{\lxNotesOuterfalse}

% =============================================================
% DIMENSIONS
%
% True units are used for all physical paper geometry.
% TeX's \magnification scales dimensions in pt/cm/mm etc. by the
% magnification ratio; "true" units (truecm, truein etc.) are exempt
% and remain physically fixed — correct for paper geometry.
%
% MARGIN WIDTH — two mutually exclusive user-facing commands:
%   \lxMarginWidth=3truecm          absolute: a true-unit dimension
%   \def\lxMarginWidthFraction{0.2} relative: a bare decimal fraction
%
% If \lxMarginWidthFraction is defined, \lxSetupPage computes
% \lxMarginWidth as that fraction of the engine-detected page width.
% (In LaTeX, \paperwidth set by documentclass — no detection needed.)
%
% ENGINE PORTABILITY — page width/height primitives:
%   pdfTeX:    \pdfpagewidth,  \pdfpageheight
%   XeTeX:     \pagewidth,     \pageheight
%   DVI/other: neither — \lxPaperWidth / \lxPaperHeight used as fallback
%              (defaults: A4 = 210truemm x 297truemm)
% \lxSetupPage detects the available primitive via \ifx...\undefined.
% The user need never mention any engine primitive directly.
% =============================================================

\newdimen\lxMarginWidth
\newdimen\lxMarginSep
\newdimen\lxInnerMargin
\newdimen\lxOuterMargin
\newdimen\lxTopMargin
\newdimen\lxBottomMargin
\newdimen\lxPaperWidth    % user-settable fallback for DVI/other engines
\newdimen\lxPaperHeight   % user-settable fallback for DVI/other engines
\newdimen\lxPageWidth     % internal: resolved at \lxSetupPage time
\newdimen\lxPageHeight    % internal: resolved at \lxSetupPage time

% Defaults:
\lxMarginWidth=3truecm    % used only if \lxMarginWidthFraction not defined
\lxMarginSep=0.5em
\lxInnerMargin=1truein
\lxOuterMargin=1truein
\lxTopMargin=1truein
\lxBottomMargin=1truein
\lxPaperWidth=210truemm   % A4; override for other paper sizes
\lxPaperHeight=297truemm  %     (not needed in LaTeX — documentclass sets this)

% \lxMarginWidthFraction intentionally NOT defined here by default.

% \lxNoteFont — complete typographic environment for margin note text.
%   Switches to \scriptfont0 and scales all skip parameters by 7/10,
%   giving a self-consistent 7pt environment at 10pt base font size.
%   (In LaTeX: \newcommand{\lxNoteFont}{\footnotesize} — approximate.)
\def\lxNoteFont{\the\scriptfont0 \spacing{7}{10}}

% =============================================================
% PAGE SETUP
%
% \lxSetupPage commits the current flag and dimension values.
% Call once after all flags are set.
%
% Sequence:
%   1. Resolve page width/height from engine primitives or fallback.
%      (In LaTeX, \paperwidth/\paperheight used — step 1 not needed.)
%   2. If \lxMarginWidthFraction defined, compute \lxMarginWidth.
%      (Same in both; in LaTeX \geometry{} follows in document preamble.)
%   3. Commit \hsize, \vsize, \hoffset, \voffset directly.
%      (In LaTeX, \geometry{} called directly in document-latex.tex.)
%   4. Activate \headline and \footline.
%      (In LaTeX, fancyhdr is wired instead — \headline not used.)
% =============================================================

\def\lxSetupPage{%
  % --- Step 1: resolve page dimensions from engine primitives ---
  \ifx\pdfpagewidth\undefined               % not pdfTeX
    \ifx\pagewidth\undefined                % not XeTeX either
      \lxPageWidth=\lxPaperWidth            % DVI/other: use fallback
      \lxPageHeight=\lxPaperHeight          %
    \else                                   %
      \lxPageWidth=\pagewidth               % XeTeX
      \lxPageHeight=\pageheight             %
    \fi                                     %
  \else                                     %
    \lxPageWidth=\pdfpagewidth              % pdfTeX
    \lxPageHeight=\pdfpageheight            %
  \fi                                       % (in LaTeX: \paperwidth/\paperheight
  %                                         %  are set by documentclass; step 1
  %                                         %  is entirely absent from lx.sty —
  %                                         %  cf. the placeholder comment block
  %                                         %  in lx.sty \lxSetupPage Step 1)
  %                                         %
  %                                         %
  %                                         %
  %                                         %
  %                                         %
  %                                         %
  %                                         %
  % --- Step 2: resolve margin width ---
  \ifx\lxMarginWidthFraction\undefined
    % absolute: \lxMarginWidth already set by user (or default)
  \else
    \lxMarginWidth=\lxMarginWidthFraction\lxPageWidth  % relative
  \fi
  % --- Step 3: commit layout registers ---
  \hsize=\lxPageWidth                       % (cf. \geometry{} in LaTeX)
    \advance\hsize by -\lxInnerMargin       %
    \advance\hsize by -\lxOuterMargin       %
    \advance\hsize by -\lxMarginWidth       %
    \advance\hsize by -\lxMarginSep         %
  \vsize=\lxPageHeight                      %
    \advance\vsize by -\lxTopMargin         %
    \advance\vsize by -\lxBottomMargin      %
  \hoffset=\lxInnerMargin \advance\hoffset by -1truein
  \voffset=\lxTopMargin   \advance\voffset by -1truein
  %
  % --- Step 4: activate lx headline and footline ---
  \headline={\lxResolvedHeadline}%          % (cf. fancyhdr wiring in LaTeX)
  \footline={\lxResolvedFootline}%          % (cf. \pagestyle{fancy}+\fancyhf{} in LaTeX)
}                                           %
%                                           (no \AtBeginDocument in Plain TeX —
%                                           \headline/\footline are not overridden
%                                           by any Plain TeX default pagestyle)
%                                           (cf. lx.sty: \AtBeginDocument{\pagestyle{fancy}})
%                                           (cf. document-latex.tex: \fancyhfoffset{0pt})

% =============================================================
% HEAD / FOOT HOOKS
%
% Defined as \newtoks registers (Plain TeX idiom). Assign with:
%   \lxHeadInner={...}
% (In LaTeX, \newcommand + \renewcommand are used instead.)
%
% Defaults reproduce Plain TeX's own \headline and \footline:
%   head: all empty (\headline={\hfil})
%   foot: \lxFootMiddle={\the\pageno} (centred page number)
% =============================================================

\newtoks\lxHeadInner  \newtoks\lxHeadMiddle  \newtoks\lxHeadOuter
\newtoks\lxFootInner  \newtoks\lxFootMiddle  \newtoks\lxFootOuter

\lxHeadInner={}  \lxHeadMiddle={}  \lxHeadOuter={}
\lxFootInner={}  \lxFootMiddle={\the\pageno}  \lxFootOuter={}

% =============================================================
% BINDING RESOLUTION
%
% \lxResolveBinding sets \iflxInnerIsLeft at output/typeset time,
% called inside \headline/\footline so \pageno is current.
% (In LaTeX, fancyhdr's [LE,RO]/[RE,LO] slot notation handles this
% automatically — no explicit resolution macro is needed.)
%
% fancyhdr slot notation (cf. \lxSetupPage in lx.sty):
%   [LE,RO] = inner (binding side, both recto and verso)
%   [RE,LO] = outer (fore-edge,    both recto and verso)
%   [L]/[R] = left/right in oneside mode
%
% Truth table (encoded in \lxResolveBinding below):
%   double, binding=left,  odd  (recto): inner=left
%   double, binding=left,  even (verso): inner=right
%   double, binding=right, odd  (recto): inner=right
%   double, binding=right, even (verso): inner=left
%   single (any binding):               inner=left always
% =============================================================

\newif\iflxInnerIsLeft

\def\lxResolveBinding{%
  \iflxDoubleSided
    \ifodd\pageno
      \iflxBindingLeft \lxInnerIsLefttrue \else \lxInnerIsLeftfalse \fi
    \else
      \iflxBindingLeft \lxInnerIsLeftfalse \else \lxInnerIsLefttrue \fi
    \fi
  \else
    \lxInnerIsLefttrue
  \fi
}

\def\lxMakeLine#1#2#3{#1\hfil#2\hfil#3}

\def\lxResolvedHeadline{%
  \lxResolveBinding
  \iflxInnerIsLeft
    \lxMakeLine{\the\lxHeadInner}{\the\lxHeadMiddle}{\the\lxHeadOuter}%
  \else
    \lxMakeLine{\the\lxHeadOuter}{\the\lxHeadMiddle}{\the\lxHeadInner}%
  \fi
}

\def\lxResolvedFootline{%
  \lxResolveBinding
  \iflxInnerIsLeft
    \lxMakeLine{\the\lxFootInner}{\the\lxFootMiddle}{\the\lxFootOuter}%
  \else
    \lxMakeLine{\the\lxFootOuter}{\the\lxFootMiddle}{\the\lxFootInner}%
  \fi
}
%                                           (\headline/\footline set in \lxSetupPage)
%                                           (cf. fancyhdr wiring in lx.sty Step 4)
%                                           (cf. lx.sty: no \iflxInnerIsLeft needed)
%                                           (cf. lx.sty: no \lxResolveBinding needed)
%                                           (cf. lx.sty: no \lxMakeLine needed)

% =============================================================
% MARGIN NOTE COUNTER AND REFERENCE FORMATTER
%
% \lxNoteNo  — count register, analogous to Plain's \pageno.
%              Incremented automatically by \lxMarginNote.
%              Starts at 0; first note displays as 1.
%              (In LaTeX: \newcounter{lxNoteNo})
%
% \lxNoteRef — formats \lxNoteNo as the inline/margin symbol.
%              Default: arabic integer via \the\lxNoteNo.
%              Override e.g.: \def\lxNoteRef{\romannumeral\lxNoteNo}
%              (In LaTeX: \renewcommand{\lxNoteRef}{\roman{lxNoteNo}})
% =============================================================

\newcount\lxNoteNo  \lxNoteNo=0            % (cf. \newcounter{lxNoteNo} in LaTeX)
%                                           (starts at 0; first note displays as 1)
%                                           (cf. lx.sty: \arabic{lxNoteNo} for display)
\def\lxNoteRef{\the\lxNoteNo}              % (cf. \newcommand{\lxNoteRef}{\arabic{lxNoteNo}})

% =============================================================
% MARGIN NOTE MACROS
%
% Two clean, unambiguous forms — no optional-argument machinery:
%
%   \lxMarginNote{NOTE}
%     Auto-incrementing form. Increments \lxNoteNo then uses
%     \lxNoteRef as the reference symbol.
%     (In LaTeX, \refstepcounter{lxNoteNo} replaces
%     \global\advance\lxNoteNo by 1.)
%
%   \lxMarginNoteRef{REF}{NOTE}
%     Explicit-reference form. Uses REF verbatim; does NOT
%     increment \lxNoteNo.
%
% Both deposit the reference symbol inline and REF + NOTE in the
% margin via \vadjust. Correct side resolved per line at typeset time.
% (In LaTeX, \vadjust is used identically; \textwidth replaces \hsize;
% \lxstyNoteRight/\lxstyNoteLeft parallel \lxNoteRight/\lxNoteLeft.)
%
% KNOWN LIMITATION: \vadjust executes during paragraph building, before
% page breaking, so \pageno is unreliable inside it. In double-sided mode
% \lxResolveBinding may read the wrong \pageno and place notes on the wrong
% side. Single-sided mode (\lxSingleSided) is unaffected. Fixing double-sided
% margin notes requires output-routine machinery not yet implemented in lx.
% (lx.sty uses the same \vadjust approach and shares this limitation.)
% =============================================================

% \lxEmitNote — shared implementation, uses \lxCurrentRef/\lxCurrentNote:
% (In LaTeX, \lxMarginNote and \lxMarginNoteRef call \vadjust directly,
% using \lxstyNoteRight/\lxstyNoteLeft in place of \lxNoteRight/\lxNoteLeft.
% The logic is identical; \textwidth replaces \hsize; \lxResolveBinding is
% replaced by \iflxNotesOuter directly since LaTeX has no \pageno-based
% binding resolution needed for single-sided mode.)
\def\lxEmitNote{%
  \raise.7ex\hbox{\the\scriptfont0 \lxCurrentRef}% % inline mark: script font at current size
  \vadjust{\vbox to 0pt{%                  % (cf. \vadjust in lx.sty — identical approach)
    \vskip-\baselineskip                   % move up to align with current line top
    \nointerlineskip                       % suppress interline glue before note content
    \lxResolveBinding                      %
    \iflxNotesOuter
      \iflxInnerIsLeft \lxNoteRight \else \lxNoteLeft \fi
    \else
      \iflxInnerIsLeft \lxNoteLeft \else \lxNoteRight \fi
    \fi
    \vss}}%
}
%                                           (\vskip-\baselineskip moves note up to current line)

\def\lxMarginNote#1{%
  \global\advance\lxNoteNo by 1            % (cf. \refstepcounter in LaTeX)
  \def\lxCurrentRef{\lxNoteRef}%           % (cf. \lxNoteRef directly in lx.sty)
  \def\lxCurrentNote{#1}%
  \lxEmitNote                              % (cf. inline \vadjust call in lx.sty)
}

\def\lxMarginNoteRef#1#2{%
  \def\lxCurrentRef{#1}%                  % (cf. #1 used directly in lx.sty)
  \def\lxCurrentNote{#2}%                 % (cf. #2 used directly in lx.sty)
  \lxEmitNote                              % (cf. inline \vadjust call in lx.sty)
}

% Deposit note content in the RIGHT margin:
% (cf. lx.sty: \lxstyNoteRight{REF}{NOTE} — identical, \textwidth for \hsize)
\def\lxNoteRight{\hbox{\kern\hsize\kern\lxMarginSep
  \vbox{\hsize=\lxMarginWidth \lxNoteFont
    \noindent\hbox{\raise.7ex\hbox{\lxCurrentRef}\thinspace}\lxCurrentNote}}}

% Deposit note content in the LEFT margin:
% (cf. lx.sty: \lxstyNoteLeft{REF}{NOTE}  — identical, \textwidth for \hsize)
\def\lxNoteLeft{\hbox{\llap{%
  \vbox{\hsize=\lxMarginWidth \lxNoteFont
    \noindent\hbox{\raise.7ex\hbox{\lxCurrentRef}\thinspace}\lxCurrentNote}%
  \kern\lxMarginSep}}}

%                                           (cf. lx.sty: \lxstyNoteRight/\lxstyNoteLeft
%                                           are named helpers taking REF and NOTE as args;
%                                           Plain TeX uses \lxCurrentRef/\lxCurrentNote
%                                           via \lxEmitNote — same net effect)

% =============================================================
% END OF lx.tex
% =============================================================
